% Vietnamese translation for OI Public Library
% Source-PDF: IOI中国国家候选队论文/2007/day2/8.陈瑜希《多角度思考 创造性思维》.pdf
\documentclass[11pt]{article}
\usepackage{oipl-vn}

\newcommand{\Oh}{\mathcal{O}}
\newcommand{\dist}{\operatorname{dist}}

\title{Suy nghĩ từ nhiều góc độ, tư duy sáng tạo}
\author{Chen Yuxi\\Trường Ngoại ngữ Nam Kinh}
\date{2007}

\begin{document}
\maketitle

\origtitle{Thinking from multiple angles and creative thinking: ideas and
methods for solving problems with tree dynamic programming}
\sourcepath{IOI China candidate team papers / 2007 / day 2 / Chen Yuxi}

\section*{Tóm tắt}

Trong vài năm gần đây, các bài toán cần dùng quy hoạch động trên cây xuất
hiện thường xuyên trong thi tin học. Những bài toán này có nhiều biến thể,
lồng ghép nhiều ý tưởng tinh tế, và đòi hỏi cao ở năng lực phân tích cũng như
tư duy sáng tạo của thí sinh. Vì vậy chúng giữ một vị trí quan trọng trong
các kỳ thi.

Bài viết phân tích một số bài toán quy hoạch động trên cây trong các kỳ thi
quốc tế và quốc gia gần đây, tập trung vào vài cách giải tiêu biểu. Từ quá
trình phân tích ấy, ta rút ra phương pháp tư duy chung: nhìn bài toán từ
nhiều góc độ và sáng tạo trong cách đặt trạng thái. Mục tiêu không chỉ là nêu
cách giải, mà còn trình bày quá trình suy nghĩ dẫn tới cách giải.

\paragraph{Từ khóa.}
Cấu trúc cây; quy hoạch động; Olympic tin học.

\section{Dẫn nhập}

Trong thi tin học thường gặp các bài toán kiểu sau: cho một cây, yêu cầu hoàn
thành một thao tác nào đó với chi phí nhỏ nhất, hoặc thu được lợi ích lớn
nhất. Rất nhiều bài toán được xây dựng bằng cách mở rộng hoặc tăng cường hai
yếu tố cơ bản là cấu trúc cây và tính tối ưu, vì vậy trở nên khá khó xử lý.

Những bài toán này thường có kích thước lớn. Thuật toán liệt kê không đủ hiệu
quả, còn tham lam thường không bảo đảm tối ưu, nên ta phải dùng quy hoạch
động.

Giống quy hoạch động nói chung, khi giải bài toán dạng này ta thường xét ba
bước.

\begin{enumerate}
  \item \term{Xác lập trạng thái.} Hầu như bài nào cũng cần lưu thông tin của
  cây con gốc tại một đỉnh nào đó. Tuy nhiên, có cần thêm chiều hay không,
  thêm bao nhiêu chiều, và thêm theo cách nào thì phải dựa vào từng bài cụ
  thể.
  \item \term{Chuyển trạng thái.} Cách chuyển rất đa dạng. Đây cũng là trọng
  tâm phân tích trong các ví dụ của bài viết.
  \item \term{Cài đặt thuật toán.} Nhờ cấu trúc cây, tìm kiếm có nhớ thường là
  cách cài đặt tự nhiên.
\end{enumerate}

Vì mô hình được xây dựng trên cây, ta gọi đó là quy hoạch động trên cây. Theo
nghĩa trực tiếp, đây là quy hoạch động trên cấu trúc dữ liệu ``cây''.

Phần lớn các bài quy hoạch động quen thuộc có cấu trúc tuyến tính. Trên đường
thẳng có hai hướng tự nhiên: tiến và lùi, tương ứng với tính xuôi và tính
ngược. Trên cây cũng có hai hướng tương tự:

\begin{enumerate}
  \item \term{Từ gốc xuống lá.} Gốc truyền thông tin hữu ích xuống các con;
  sau đó nghiệm tối ưu được suy ra từ thông tin đã truyền.
  \item \term{Từ lá lên gốc.} Các con truyền thông tin hữu ích lên cha; sau
  đó gốc nhận được nghiệm tối ưu. Đây là dạng thường gặp hơn và có phạm vi
  ứng dụng rộng.
\end{enumerate}

Dĩ nhiên cũng có những bài cần đồng thời cả hai hướng duyệt. Ví dụ đầu tiên
dưới đây thuộc loại đó.

\section{Ví dụ 1: Tìm Chris}

\subsection{Mô tả bài toán}

Điện thoại nhà Chris reo lên. Từ đầu dây bên kia vang lên giọng hốt hoảng của
giáo viên: ``A lô, có phải phụ huynh của Chris không? Con anh chị lại không
đến lớp, không định tham gia kỳ thi sao?'' Nghe đến chuyện thi, bố mẹ Chris
rất lo và quyết định tìm Chris trong thời gian ngắn nhất. Họ nói với giáo
viên: ``Theo kinh nghiệm trước đây, bây giờ Chris chắc chắn đang trốn ở nhà
bạn Shermie hoặc Yashiro để chơi \emph{The King of Fighters}. Chúng tôi sẽ đi
từ nhà để tìm Chris; hễ tìm thấy sẽ gọi lại ngay.'' Nói xong họ cúp máy.

Thành phố nơi Chris sống gồm \(N\) điểm dân cư và một số đường hai chiều nối
giữa chúng. Đi qua đường \(x\) mất \(T_x\) phút. Giữa hai điểm dân cư bất kỳ
có đúng một đường đi, tức mạng đường là một cây. Nhà Chris ở điểm \(C\),
Shermie ở điểm \(A\), còn Yashiro ở điểm \(B\). Giáo viên của Chris và bố mẹ
Chris đều có bản đồ thành phố; bố mẹ biết vị trí cụ thể của \(A,B,C\), còn
giáo viên thì không.

Để nhanh chóng tìm Chris, bố mẹ cậu tuân theo hai quy tắc:
\begin{itemize}
  \item Nếu \(A\) gần \(C\) hơn \(B\) gần \(C\), họ đến nhà Shermie trước; nếu
  không thấy Chris, họ đến nhà Yashiro. Trường hợp ngược lại xử lý tương tự.
  \item Họ luôn đi theo đường đi duy nhất giữa hai điểm.
\end{itemize}

Giáo viên biết hai quy tắc trên, nhưng không biết vị trí cụ thể của
\(A,B,C\). Hãy cho biết trong trường hợp xấu nhất, bố mẹ Chris phải mất bao
lâu để tìm được Chris.

Ví dụ, xét cây bốn đỉnh sau, mọi cạnh có độ dài \(1\).

\begin{verbatim}
    A -- C -- B
         |
       Chris
\end{verbatim}

Nếu Chris sống ở \(C\), Shermie ở \(A\), Yashiro ở \(B\), và \(B\) gần \(C\)
hơn \(A\), bố mẹ Chris sẽ đến nhà Yashiro trước. Nếu không thấy Chris ở đó,
họ đi tiếp đến nhà Shermie. Khi đó trường hợp xấu nhất mất \(4\) phút.

\paragraph{Dữ liệu vào và ra.}
Tệp \texttt{hookey.in} có dòng đầu gồm hai số nguyên \(N\) và \(M\)
\((3\le N\le 200000)\), lần lượt là số điểm dân cư và số đường. Mỗi dòng
trong \(M\) dòng tiếp theo chứa \(U_i,V_i,T_i\)
\((1\le U_i,V_i\le N,\ 1\le T_i\le 10^9)\), nghĩa là đường \(i\) nối
\(U_i\) với \(V_i\) và đi qua nó mất \(T_i\) phút. Mỗi đường được cho đúng
một lần. Tệp \texttt{hookey.out} chứa một số nguyên \(T\), là thời gian lớn
nhất cần thiết trong trường hợp xấu nhất.

\subsection{Phân tích}

\paragraph{Trừu tượng hóa.}
Bài toán tương đương với việc chọn ba đỉnh \(X,Y,Z\) trên một cây có trọng số
sao cho
\[
  \dist(X,Y)\le \dist(X,Z),
\]
và giá trị
\[
  \dist(X,Y)+\dist(Y,Z)
\]
lớn nhất có thể. Ở đây \(X\) là nhà Chris, \(Y\) là nơi được đến trước, còn
\(Z\) là nơi được đến sau.

\paragraph{Phân tích ban đầu.}
Mô hình là một cây và kích thước rất lớn, nên ta dễ nghĩ đến quy hoạch động
trên cây. Nếu cố định một đỉnh \(Y\) làm nơi đến trước, ta chỉ cần tìm hai
đỉnh xa nhất liên quan tới \(Y\). Nhưng trên cây, tìm hai đỉnh xa nhất từ một
đỉnh bất kỳ tốn \(\Oh(n)\); vì chưa biết đỉnh nào là \(Y\), nếu liệt kê
\(Y\) thì tổng thời gian thành \(\Oh(n^2)\), không thể chạy với
\(N=200000\).

Không liệt kê trực tiếp \(Y\) không có nghĩa là phải bỏ hẳn ý tưởng liệt kê.
Ta thử nhìn bài toán từ một góc khác. Xét dạng đường đi sau:

\begin{verbatim}
        Y
        |
        a
       / \
      X   Z
\end{verbatim}

Hai đường đi \(Y\to X\) và \(Y\to Z\) cùng đi từ \(Y\) tới một đỉnh \(a\) rồi
tách ra. Ta gọi \(a\) là \term{điểm phân nhánh}. Nếu liệt kê điểm phân nhánh,
liệu có thể có thuật toán hiệu quả hơn không?

\paragraph{Liệt kê điểm phân nhánh.}
Giả sử chọn một đỉnh \(a\) làm điểm phân nhánh và gốc hóa cây tại \(a\). Đỉnh
\(Y\) có hai khả năng: \(Y=a\), hoặc \(Y\) nằm trong một cây con của một con
của \(a\). Trường hợp \(Y=a\) có thể xem như trường hợp thứ hai bằng cách thêm
cho \(a\) một con rỗng có cạnh dài \(0\).

Vì \(a\) là điểm phân nhánh, ba đỉnh \(X,Y,Z\) phải nằm ở ba nhánh khác nhau
đi ra từ \(a\). Khi đó
\[
  \begin{aligned}
  \dist(X,Y)+\dist(Y,Z)
  &=(\dist(X,a)+\dist(Y,a))+(\dist(Y,a)+\dist(Z,a))\\
  &=\dist(X,a)+2\dist(Y,a)+\dist(Z,a).
  \end{aligned}
\]
Điều kiện \(\dist(X,Y)\le\dist(X,Z)\) tương đương
\(\dist(Y,a)\le\dist(Z,a)\). Vì vậy, nếu biết ba khoảng cách lớn nhất đi từ
\(a\) vào ba nhánh khác nhau và sắp chúng thành
\(d_1\ge d_2\ge d_3\), giá trị tốt nhất tại \(a\) là
\[
  d_1+2d_2+d_3.
\]

Nếu với mỗi \(a\) ta lại quét toàn cây để tìm ba giá trị này, thuật toán vẫn
là \(\Oh(n^2)\). Ta cần một cách truyền thông tin tốt hơn.

\paragraph{Hai lần duyệt.}
Gốc hóa cây tại đỉnh \(1\). Lần duyệt thứ nhất đi từ lá lên gốc để tính, với
mỗi đỉnh, các khoảng cách xa nhất nằm trong những cây con của nó. Lần duyệt
thứ hai đi từ gốc xuống lá, chẳng hạn theo BFS để cha luôn được xử lý trước
con. Khi đang xét đỉnh \(a\), một điểm xa nhất từ \(a\) hoặc nằm trong cây con
của \(a\), hoặc nằm ngoài cây con đó. Nếu nằm ngoài, đường đi chắc chắn phải
qua cha của \(a\).

Do đó tại mỗi đỉnh ta lưu hai giá trị lớn nhất được truyền tới nó, đồng thời
lưu mỗi giá trị đến từ đỉnh kề nào. Khi truyền thông tin từ một đỉnh sang một
con \(v\), nếu giá trị lớn nhất của đỉnh hiện tại đến từ chính \(v\) thì dùng
giá trị lớn thứ hai; nếu không thì dùng giá trị lớn nhất. Cộng thêm độ dài
cạnh giữa hai đỉnh để được khoảng cách nhìn từ \(v\).

Sau hai lần duyệt, với mỗi đỉnh ta đã có các ứng viên xa nhất ở những nhánh
khác nhau. Lấy ba giá trị lớn nhất thuộc ba nguồn kề khác nhau, sắp giảm dần
thành \(d_1,d_2,d_3\), rồi cập nhật đáp án bằng \(d_1+2d_2+d_3\). Cả thời gian
và bộ nhớ đều là \(\Oh(n)\).

\paragraph{Cài đặt.}
Các giá trị trung gian và đáp án có thể vượt phạm vi số nguyên \(32\) bit, vì
vậy cần dùng kiểu \(64\) bit.

Ý tưởng trong ví dụ này có tính sáng tạo: thay vì nhìn vào đỉnh bắt đầu hay
đỉnh đến trước, ta nhìn vào điểm phân nhánh của hai đường đi. Phương pháp hai
lần duyệt và kỹ thuật ``lưu vài giá trị lớn nhất cùng nguồn truyền'' xuất
hiện trong rất nhiều bài toán trên cây.

\section{Ví dụ 2: Nhà máy gỗ ở Byteland}

\subsection{Mô tả bài toán}

Gần như toàn bộ vương quốc Byteland được bao phủ bởi rừng và sông. Các sông
nhỏ hợp lại thành sông lớn hơn; cuối cùng mọi dòng nước đều đổ vào một con
sông lớn, rồi con sông này chảy ra biển. Tại cửa sông có một ngôi làng tên là
Bytetown.

Trong Byteland có \(n\) làng khai thác gỗ, đều nằm ven sông. Hiện tại ở
Bytetown có một nhà máy gỗ lớn xử lý toàn bộ gỗ của cả nước. Sau khi được
chặt, gỗ trôi theo sông về nhà máy ở Bytetown. Nhà vua quyết định xây thêm
\(k\) nhà máy gỗ ở các làng khác để giảm chi phí vận chuyển. Khi đó gỗ sẽ
được xử lý tại nhà máy mới đầu tiên mà nó gặp trên đường trôi xuống hạ lưu.
Nếu nhà máy đặt ngay tại làng sinh ra gỗ thì gỗ không tốn chi phí vận chuyển.

Tất cả các dòng sông đều không phân nhánh theo chiều xuôi dòng: từ mỗi làng,
đi theo hạ lưu chỉ có một đường đến Bytetown. Mỗi làng có sản lượng gỗ hằng
năm đã biết. Chi phí vận chuyển là \(1\) xu cho mỗi khúc gỗ trên mỗi kilomet.
Hãy chọn nơi xây \(k\) nhà máy mới để tổng chi phí nhỏ nhất.

\paragraph{Dữ liệu vào và ra.}
Dòng đầu gồm \(n\) và \(k\)
\((2\le n\le 100,\ 1\le k\le 50,\ k\le n)\). Các làng ngoài Bytetown được
đánh số \(1,2,\ldots,n\), còn Bytetown là \(0\). Trong \(n\) dòng tiếp theo,
dòng của làng \(i\) gồm:
\[
  w_i,\quad v_i,\quad d_i,
\]
trong đó \(w_i\) là số khúc gỗ làng \(i\) sinh ra mỗi năm
\((0\le w_i\le 10000)\), \(v_i\) là làng gần nhất ở hạ lưu của \(i\)
\((0\le v_i\le n)\), và \(d_i\) là khoảng cách từ \(v_i\) đến \(i\)
\((1\le d_i\le 10000)\). Kết quả là chi phí vận chuyển nhỏ nhất. Tổng chi phí
nếu đưa toàn bộ gỗ về Bytetown không vượt quá \(2\,000\,000\,000\) xu.

\subsection{Phân tích}

\paragraph{Trừu tượng hóa.}
Ta có một cây có gốc là Bytetown. Cần đặt \(k\) nhà máy ở các đỉnh khác gốc,
sao cho mỗi làng đưa gỗ tới nhà máy tổ tiên gần nhất và tổng chi phí nhỏ
nhất. Đây là một bài toán quy hoạch động trên cây khá tự nhiên.

\paragraph{Xác lập trạng thái.}
Trước hết trạng thái phải biết đỉnh hiện tại và số nhà máy được đặt trong cây
con gốc tại đỉnh đó. Nhưng như vậy chưa đủ, vì chi phí còn phụ thuộc vào vị
trí nhà máy tổ tiên gần nhất ở phía trên. Ta thêm một chiều biểu diễn tổ tiên
đang được dùng làm nơi nhận gỗ nếu trong phần đang xét chưa gặp nhà máy mới.

Đặt
\[
  f(u,t,q)
\]
là chi phí nhỏ nhất trong cây con gốc tại \(u\), khi đặt \(q\) nhà máy trong
cây con này, và \(t\) là một tổ tiên của \(u\) đã có nhà máy. Ở đây \(t\) chỉ
cần là một nhà máy tổ tiên đang được truyền xuống, không nhất thiết phải mô tả
mọi chi tiết của đường đi phía trên. Chính việc nới rộng cách hiểu trạng thái
này làm chuyển trạng thái đơn giản hơn.

\paragraph{Chuyển trạng thái.}
Xét hai trường hợp.

\term{Đặt nhà máy tại \(u\).} Khi đó gỗ của \(u\) không tốn chi phí vận
chuyển, và các con của \(u\) xem \(u\) là nhà máy tổ tiên. Một nhà máy đã được
dùng tại \(u\), nên số nhà máy còn lại được phân phối cho các cây con. Quá
trình phân phối này giống bài toán ba lô trên các con:
\[
  h_1(i,s)=
  \min_{0\le r\le s}
  \bigl\{h_1(i-1,s-r)+f(\operatorname{son}_i,u,r)\bigr\}.
\]

\term{Không đặt nhà máy tại \(u\).} Khi đó gỗ của \(u\) phải đi tới \(t\), tạo
chi phí \(w_u\dist(u,t)\). Các con vẫn dùng \(t\) làm nhà máy tổ tiên đang
được truyền xuống:
\[
  h_2(i,s)=
  \min_{0\le r\le s}
  \bigl\{h_2(i-1,s-r)+f(\operatorname{son}_i,t,r)\bigr\}.
\]

Vì vậy, nếu \(u\) có \(c\) con,
\[
  f(u,t,q)=
  \min\left\{
    h_1(c,q-1),
    w_u\dist(u,t)+h_2(c,q)
  \right\}.
\]
Các mảng \(h_1,h_2\) là mảng tạm để gộp dần các con. Trường hợp cơ sở được
xử lý theo số con bằng \(0\) và số nhà máy còn lại hợp lệ.

\paragraph{Hiệu quả.}
Trạng thái có ba chiều. Trong chuyển trạng thái ta phân phối số nhà máy giữa
các con, tương tự ba lô. Với \(n\le 100\) và \(k\le 50\), cách cài đặt chuẩn
có thể đạt cỡ \(\Oh(n^3k)\), đủ cho ràng buộc. Điểm đáng chú ý của bài là
phải thêm chiều cho vị trí nhà máy tổ tiên, và trong phạm vi các anh em lại
thực hiện một quy hoạch động thứ hai kiểu ba lô.

\section{Ví dụ 3: Tính phí mạng theo cặp}

\subsection{Mô tả bài toán}

Mạng đã trở thành một phần không thể thiếu của thế giới hiện nay. Mỗi ngày có
hàng trăm triệu người dùng mạng để học tập, nghiên cứu, giải trí. Tuy nhiên,
bản thân mạng có chi phí vận hành rất lớn; vì vậy thu phí người dùng ở mức
thích hợp là cần thiết và hợp lý.

Trường NS ở thành phố MY có một mạng giáo dục như sau. Mạng có \(2^N\) người
dùng, đánh số \(1,2,\ldots,2^N\). Người dùng, điểm định tuyến và dây mạng tạo
thành một cây nhị phân đầy đủ. Mỗi lá là một người dùng, mỗi đỉnh trong là một
điểm định tuyến, mỗi cạnh là một dây mạng. Ví dụ với \(N=3\):

\begin{verbatim}
              *
          /       \
        *           *
      /   \       /   \
     *     *     *     *
    / \   / \   / \   / \
   1   2 3   4 5   6 7   8
\end{verbatim}

Công ty mạng MY có cách thu phí đặc biệt gọi là ``thu phí theo cặp''. Với mỗi
cặp người dùng \(i,j\) \((1\le i<j\le 2^N)\), công ty tính một khoản phí. Mỗi
người dùng tự chọn một trong hai phương thức trả phí \(A\) hoặc \(B\). Tổng
phí mà trường phải trả là tổng đóng góp của mọi cặp.

Với hai người dùng \(i,j\), gọi \(P\) là tổ tiên chung gần nhất của họ. Trong
các lá thuộc quyền quản lý của \(P\), gọi \(n_A,n_B\) lần lượt là số người
chọn \(A\) và \(B\). Hệ số tính phí \(k\) cho cặp \(i,j\) được cho bởi bảng
sau; \(F_{i,j}\) là lưu lượng giữa hai người dùng này.

\begin{center}
\begin{tabular}{c c c c}
\hline
Phương thức của \(i\) & Phương thức của \(j\) & Điều kiện & Hệ số \(k\)\\
\hline
\(A\) & \(A\) & \(n_A<n_B\) & \(2\)\\
\(A\) & \(B\) & \(n_A<n_B\) & \(1\)\\
\(B\) & \(A\) & \(n_A<n_B\) & \(1\)\\
\(B\) & \(B\) & \(n_A<n_B\) & \(0\)\\
\hline
\(A\) & \(A\) & \(n_A\ge n_B\) & \(0\)\\
\(A\) & \(B\) & \(n_A\ge n_B\) & \(1\)\\
\(B\) & \(A\) & \(n_A\ge n_B\) & \(1\)\\
\(B\) & \(B\) & \(n_A\ge n_B\) & \(2\)\\
\hline
\end{tabular}
\end{center}

Phí thực tế của cặp là \(kF_{i,j}\). Vì tổng phí phụ thuộc vào lựa chọn
phương thức, học sinh muốn thay đổi phương thức để giảm phí. Nhưng công ty đã
lưu phương thức đăng ký ban đầu; nếu người dùng \(i\) muốn đổi từ \(A\) sang
\(B\) hoặc ngược lại thì phải trả thêm \(C_i\) để sửa hồ sơ. Hãy tìm tổng chi
phí nhỏ nhất, gồm phí sửa hồ sơ và phí theo cặp.

\paragraph{Dữ liệu vào và ra.}
Dòng đầu là số nguyên dương \(N\). Dòng thứ hai có \(2^N\) số, mô tả phương
thức đăng ký ban đầu của từng người dùng; \(0\) nghĩa là \(A\), \(1\) nghĩa
là \(B\). Dòng thứ ba có \(2^N\) số \(C_1,C_2,\ldots,C_{2^N}\). Các dòng còn
lại cho bảng lưu lượng \(F\): ở dòng tương ứng với \(i\), cột \(j\) cho
\(F_{i,i+j}\). Cần in ra tổng chi phí nhỏ nhất.

Ràng buộc: \(40\%\) dữ liệu có \(N\le 4\), \(80\%\) dữ liệu có \(N\le 7\),
toàn bộ dữ liệu có \(N\le 10\), \(0\le F_{i,j}\le 500\),
\(0\le C_i\le 500000\).

\subsection{Phân tích}

\paragraph{Biến đổi bài toán.}
Ràng buộc gợi ý quy hoạch động trên cây, nhưng quy tắc thu phí theo cặp làm
trạng thái ban đầu khó đặt. Quan sát bảng hệ số, ta thấy nó có dạng rất đặc
biệt:
\[
  2,1,1,0
  \quad\text{và}\quad
  0,1,1,2.
\]
Với hai người dùng \(i,j\), gọi LCA của họ là \(p\). Khi \(n_A<n_B\) tại
\(p\), mỗi đầu mút chọn \(A\) đóng một lần \(F_{i,j}\). Khi \(n_A\ge n_B\),
mỗi đầu mút chọn \(B\) đóng một lần \(F_{i,j}\). Nhờ vậy, phí theo cặp có thể
được chuyển thành phí gán cho từng người dùng trong cặp, tổng cuối cùng không
đổi. Đây là tiền đề để làm quy hoạch động.

\paragraph{Xác lập trạng thái.}
Trạng thái cần biết trong miền của một đỉnh \(i\) có bao nhiêu người dùng
chọn \(A\). Nhưng như vậy vẫn chưa đủ, vì một người dùng còn bị thu phí bởi
các cặp có LCA ở các tổ tiên của \(i\); tại mỗi tổ tiên ấy, việc thu theo
\(A\) hay \(B\) phụ thuộc vào so sánh \(n_A,n_B\).

Đặt
\[
  f(i,j,k)
\]
là chi phí nhỏ nhất trong miền lá do đỉnh \(i\) quản lý, khi có đúng \(j\)
người chọn \(A\). Với mỗi tổ tiên \(u\) của \(i\), ghi nhãn \(0\) nếu tại
\(u\) có \(n_A<n_B\), và nhãn \(1\) nếu \(n_A\ge n_B\). Dãy nhãn của các tổ
tiên được mã hóa bằng số nhị phân \(k\). Trạng thái \(f(i,j,k)\) lưu chi phí
nhỏ nhất dưới điều kiện đó.

\paragraph{Chuyển trạng thái.}
Giả sử \(i\) có con trái \(L\) và con phải \(R\). Ta phân phối \(j\) người
chọn \(A\) giữa hai con. Nếu con trái có \(u\) người chọn \(A\), con phải có
\(j-u\). Từ \(j\) và kích thước miền của \(i\), ta biết nhãn mới \(s\) của
đỉnh \(i\):
\[
  s=
  \begin{cases}
    0, & j < \operatorname{size}(i)-j,\\
    1, & j \ge \operatorname{size}(i)-j.
  \end{cases}
\]
Nhãn này được thêm vào dãy tổ tiên khi đi xuống hai con, tức mã mới là
\(\operatorname{push}(k,s)\). Khi đó
\[
  f(i,j,k)=
  \min_u
  \left\{
    f(L,u,\operatorname{push}(k,s))
    +f(R,j-u,\operatorname{push}(k,s))
  \right\}.
\]
Các khoản phí phát sinh giữa hai miền con được tính ở mức \(i\), vì LCA của
mọi cặp một lá bên trái và một lá bên phải chính là \(i\). Cách viết trên gom
phần ấy vào chi phí con hoặc cộng thêm khi hợp nhất, tùy cách cài đặt.

\paragraph{Biên.}
Khi \(i\) là lá, từ dãy nhãn \(k\) ta biết người dùng ở lá này phải đóng phí
cho các cặp qua những tổ tiên nào theo phương thức \(A\) hay \(B\). Cộng thêm
chi phí đổi phương thức nếu phương thức cuối cùng khác phương thức đăng ký ban
đầu.

\paragraph{Hiệu quả.}
Đặt \(m=2^N\). Nhìn thô, không gian có vẻ là \(\Oh(m^3)\) và thời gian là
\(\Oh(m^4)\), vượt xa giới hạn. Tuy nhiên rất nhiều trạng thái không thể xảy
ra.

Xem gốc là tầng \(0\), lá là tầng \(N\). Ở tầng \(i\), mỗi đỉnh quản lý nhiều
nhất \(2^{N-i}\) lá, nên số giá trị khả dĩ của \(j\) chỉ là \(2^{N-i}+1\).
Số tổ tiên chỉ là \(i\), nên số mã \(k\) nhiều nhất là \(2^i\). Tích hai chiều
sau của mỗi đỉnh là \(\Oh(2^{N-i}2^i)=\Oh(m)\). Nếu tổ chức bộ nhớ theo tầng
hoặc theo tìm kiếm có nhớ, không gian thực tế là \(\Oh(m^2)\).

Về thời gian, tại tầng \(i\) có \(2^i\) đỉnh. Mỗi đỉnh có \(\Oh(m)\) trạng
thái, và mỗi chuyển trạng thái phân phối số lượng \(A\) trong phạm vi
\(\Oh(2^{N-i})\). Tổng thời gian mỗi tầng là \(\Oh(m^2)\), nên toàn bộ thời
gian là
\[
  \Oh(m^2N).
\]
Cài đặt nên dùng tìm kiếm có nhớ và thao tác bit để xử lý mã \(k\).

Ngoài thiết kế quy hoạch động, bài này có hai điểm khó: biến đổi quy tắc thu
phí, và phân tích lại độ phức tạp bằng cách đếm trạng thái khả dĩ. Trong các
bài trên cây, nhiều thuật toán nhìn qua có vẻ quá chậm, nhưng sau khi phân
tích phân bổ theo tầng hoặc theo kích thước cây con, độ phức tạp thực tế có
thể thấp hơn một chiều hoặc hơn nữa.

\section{Ví dụ 4: Hang động và trò đoán kho báu}

\subsection{Mô tả bài toán}

Nước S có một hang động gồm \(n\) phòng và một số hành lang. Giữa hai phòng
bất kỳ luôn có đúng một đường đi. A giấu nhiều kho báu trong một phòng nhưng
không nói cho B biết phòng nào. B muốn tìm vị trí kho báu nên được đoán nhiều
lần; mỗi lần B đoán một phòng có kho báu. Nếu đoán đúng, A nói rằng B đã
đúng. Nếu đoán sai, A cho biết từ phòng B vừa đoán phải đi qua hành lang nào
đầu tiên để tới kho báu; A chỉ nói hành lang đầu tiên đó.

Hãy đọc mô tả hang động từ \texttt{jas.in}, tính số lần đoán ít nhất cần thiết
trong trường hợp xấu nhất để B biết vị trí kho báu, rồi ghi kết quả ra
\texttt{jas.out}.

\paragraph{Dữ liệu vào và ra.}
Dòng đầu có \(n\) \((1\le n\le 50000)\). Các phòng đánh số từ \(1\) đến
\(n\). Trong \(n-1\) dòng tiếp theo, mỗi dòng chứa hai số \(a,b\)
\((1\le a,b\le n)\), nghĩa là có một hành lang giữa \(a\) và \(b\). Chương
trình in ra một số nguyên: số lần đoán ít nhất cần thiết trong trường hợp xấu
nhất.

\paragraph{Ví dụ.}
\begin{verbatim}
jas.in
5
1 2
2 3
4 3
5 3

jas.out
2
\end{verbatim}

\subsection{Phân tích}

Trước hết, ta dễ đoán rằng bài toán liên quan tới quy hoạch động trên cây.
Nhưng nếu chỉ đặt trạng thái là ``với cây con gốc tại \(i\), cần ít nhất bao
nhiêu lần hỏi'' thì chưa đủ.

\paragraph{Xác lập trạng thái.}
Sau mỗi lần hỏi, B được biết kho báu nằm trong một thành phần liên thông nào
đó. Nếu thành phần còn lại không chứa đỉnh \(i\), thì số lần hỏi đã dự tính
cho cây con gốc \(i\) là đủ. Ngược lại, nếu thành phần vẫn chứa \(i\), việc
hỏi có thể đã lãng phí một phần thông tin: kho báu còn nằm trong cây con,
nhưng cấu hình chưa chắc giống như sau khi giảm đúng một lần hỏi.

Vì vậy ta cần một dãy tham số. Với cây con gốc tại \(i\), đặt
\[
  Z_{i,0},Z_{i,1},Z_{i,2},\ldots
\]
để mô tả các mức ``dư địa'' sau những lần hỏi mà thành phần còn lại vẫn chứa
\(i\). Trực giác là: nếu sau lần hỏi đầu tiên ta vẫn biết kho báu ở trong cây
con và thành phần chứa \(i\), thì cần ít nhất \(Z_{i,1}\) lần hỏi nữa trong
tình huống xấu nhất; sau lần thứ hai tương tự là \(Z_{i,2}\), v.v. Khi chọn
chiến lược, ta muốn \(Z_{i,0}\) nhỏ nhất; nếu hòa thì muốn \(Z_{i,1}\) nhỏ
nhất; tiếp tục như vậy.

Điểm đột phá là không chỉ hỏi ``cây con này cần bao nhiêu lần'', mà còn ghi
lại nếu các lần hỏi trước vẫn chưa loại được phía chứa gốc thì ta còn cần bao
nhiêu. Có những trường hợp hai phương án cùng cần số lần hỏi tối đa như nhau,
nhưng phương án có dãy \(Z\) tốt hơn sẽ hữu ích hơn khi ghép với các cây con
khác.

\paragraph{Hai ví dụ chuyển trạng thái.}
Xét một đỉnh \(A\) có hai con \(i\) và \(j\).

Ví dụ thứ nhất:
\[
  Z_i=(5,3,1),\qquad Z_j=(4,2).
\]
Ta có thể hỏi trước trong cây con \(i\). Nếu kho báu nằm trong \(i\), \(5\)
lần hỏi là đủ. Nếu không, chuyển sang cây con \(j\); nếu kho báu ở đó, cần
\(4\) lần nữa, cộng với một lần trước là \(5\). Nếu vẫn không rơi vào nhánh
đó, quay lại xét mức tiếp theo của \(i\), cần \(3\) lần nữa cộng hai lần
trước, cũng là \(5\). Vì vậy tổng cộng \(5\) lần là đủ mà không cần hỏi ngay
tại \(A\).

Ví dụ thứ hai:
\[
  Z_i=(4,2,1),\qquad Z_j=(2).
\]
Nếu hỏi trong cây con \(i\) trước và kho báu nằm trong đó, cần \(4\) lần. Nếu
không, ta hỏi một lần tại \(A\), khi đó biết chắc kho báu ở nhánh \(i\) hay
nhánh \(j\). Dù ở nhánh nào, số lần còn lại đều là \(2\), nên tổng cộng cũng
là \(4\). Ở đây có lúc phải hỏi tại \(A\), nhưng không nhất thiết hỏi \(A\)
ngay từ đầu.

\paragraph{Cài đặt thuật toán.}
Với mỗi đỉnh \(i\), chương trình lưu một bảng \(Z_i\) có độ dài
\(\lceil\log_2 n\rceil+O(1)\); có thể chứng minh số lần hỏi cỡ logarit là đủ.
Với lá, hiển nhiên
\[
  Z_{i,0}=1,
\]
các vị trí còn lại bằng \(0\).

Với đỉnh hiện tại \(A\), trước hết cộng các dãy \(Z\) của mọi con theo từng
vị trí:
\[
  Z_A=\sum_{\text{con }v\text{ của }A} Z_v.
\]
Nếu một vị trí của \(Z_A\) có giá trị \(1\), ta hiểu rằng có đúng một cây con
cần được hỏi ở mức đó; ta có thể hỏi trong cây con ấy theo cách của ví dụ
thứ nhất.

Gọi \(x\) là vị trí lớn nhất sao cho \(Z_{A,x}>1\). Ta cần tìm vị trí nhỏ
nhất \(y>x\) với \(Z_{A,y}=0\). Khi đó hỏi một lần tại \(A\) ở mức \(y\), đặt
\[
  Z_{A,y}=1,\qquad Z_{A,0}=Z_{A,1}=\cdots=Z_{A,y-1}=0.
\]
Nếu lần hỏi tại \(A\) trả lời rằng kho báu nằm trên đường đi về cha, thì chắc
chắn kho báu không nằm trong cây con gốc \(A\), nên không cần các mức thấp
hơn nữa. Đây chính là dãy \(Z\) của đỉnh \(A\).

Vì sao khi hòa phải chọn dãy \(Z\) nhỏ hơn theo thứ tự từ thấp lên cao? Hãy
xem dãy \(Z\) như một số nhị phân, trong đó \(Z_0\) là bit thấp nhất. Khi đó
\((5,2)\) tương ứng với \(100100_2\), còn \((4,3)\) tương ứng với
\(011000_2\); dãy sau nhỏ hơn. Nếu một dãy lớn hơn vừa khớp lệch với dãy của
cây con khác và tạo thành nhiều bit \(1\), về sau có thể buộc phải hỏi ở gốc
sớm hơn. Chọn dãy nhỏ nhất ngay từ đầu cho kết quả tốt hơn.

Bài này khó hơn ba ví dụ trước vì trạng thái không hiện ra trực tiếp. Ta phải
bắt đầu từ các ví dụ nhỏ, cảm nhận điều gì đang bị lãng phí sau mỗi lần hỏi,
rồi mới tìm được dãy trạng thái phù hợp.

\section{Tổng kết}

Bốn ví dụ trên minh họa nhiều phương pháp giải bài toán quy hoạch động trên
cây: duyệt cây hai lần; ý tưởng điểm phân nhánh; lưu vài giá trị tốt nhất và
nguồn truyền của chúng; quy hoạch động nhiều chiều; làm thêm một quy hoạch
động kiểu ba lô giữa các con; và phân tích độ phức tạp theo cách phân bổ trên
cấu trúc cây. Những kỹ thuật này đều có phạm vi ứng dụng rộng trong thi đấu.

Qua các ví dụ, ta thấy khi gặp một bài toán lạ, một mặt cần phân tích sâu các
thuộc tính của bài toán để tìm bản chất, mặt khác cần nhìn từ nhiều góc độ và
nắm lấy tính chất đặc biệt của nó. Đó là con đường để tạo ra cách giải đúng.

Tất nhiên, các bài toán kiểu này có rất nhiều biến thể. Bài viết chỉ đưa ra
một số phương pháp và hướng suy nghĩ; điều quan trọng vẫn là tích lũy và tổng
kết thường xuyên trong quá trình luyện tập.

\begin{thebibliography}{9}
\bibitem{practical}
Wu Wenhu, Wang Jiande. \emph{Analysis and Program Design of Practical
Algorithms}. Publishing House of Electronics Industry.

\bibitem{algds}
Fu Qingxiang, Wang Xiaodong. \emph{Algorithms and Data Structures}. Publishing
House of Electronics Industry.

\bibitem{datastructures}
Yan Weimin, Wu Weimin. \emph{Data Structures}, second edition. Tsinghua
University Press.

\bibitem{clrs}
Thomas H. Cormen, Charles E. Leiserson, Ronald L. Rivest, Clifford Stein.
\emph{Introduction to Algorithms}. Higher Education Press and The MIT Press.
\end{thebibliography}

\end{document}
